\documentclass[11pt]{article}

\usepackage[review]{acl}

\usepackage{times}
\usepackage{latexsym}
\usepackage[T1]{fontenc}
\usepackage[utf8]{inputenc}
\usepackage{microtype}
\usepackage{inconsolata}
\usepackage{graphicx}
\usepackage{booktabs}
\usepackage{multirow}

\title{Cross-Lingual Sentiment Disparity toward Geopolitical Adversaries}

\author{}

\begin{document}
\makeatletter
\acl@anonymizefalse
\makeatother
\maketitle

\begin{abstract}
Multilingual sentiment classifiers are trained once and applied across many languages, yet the same country can receive different sentiment predictions depending on prompt language. We test whether XLM-RoBERTa-Twitter assigns more negative sentiment to U.S.-designated adversary countries in English than in Chinese, Russian, or Arabic. Using 30 neutral synthetic templates and 30 countries balanced across adversary, ally, and neutral categories, we define a per-country Cross-Lingual Sentiment Gap as English bipolar valence minus the mean of the three non-English scores. A one-way ANOVA on Gap is significant (F(2, 27) = 8.75, p = .001, $\eta^2$ = .39), and a directional contrast confirms adversaries have a more negative English gap than allies and neutrals combined (t(27) = $-$4.16, one-tailed p = .0001, Cohen's d = $-$1.61 between adversary and non-adversary groups). Ally-versus-neutral is null. A nonsense-country baseline produces Gaps of comparable magnitude to adversaries, placing adversary names in the same descriptive range as country-identity-free input while ally and neutral names sit above it. This is consistent with the model encoding geopolitically-aligned sentiment differences that track the framing of its English-heavy training distribution, although the baseline is descriptive and does not by itself localize the asymmetry to a specific language side.
\end{abstract}

\section{Introduction}

A multilingual sentiment classifier is meant to behave consistently across the languages it covers: translating the same input into a different language should not change its predicted sentiment. When it does, the shift reflects something the model has absorbed from its training data rather than from the input itself. This matters because these models are used to measure public opinion, moderate content, and rank search results, and a systematic tilt toward or against a country in one language will propagate into those pipelines.

We test whether countries labeled as geopolitical adversaries in the U.S. framework receive more negative sentiment predictions in English than in three other major world languages (Chinese, Russian, Arabic). The prediction is that this English-versus-non-English gap will be larger in magnitude for adversary nations than for neutral or allied ones, on the reasoning that English-heavy training data over-represents U.S.-aligned discourse.

Our contribution is a controlled, template-based diagnostic for this question. We build 30 neutral English templates across five structural families, translate them into the three non-English languages with fluent-speaker or back-translation review, and substitute 30 countries balanced across adversary, ally, and neutral categories. We define the Cross-Lingual Sentiment Gap as a single per-country scalar and test three category means against it with a one-way ANOVA and a directional contrast. A nonsense-country baseline is run in parallel to check that any observed effect is country-identity-driven rather than a template artifact.

Two features differ from prior work. Existing studies on bias in multilingual sentiment analysis focus on demographic groups or on cross-lingual transfer effects in aggregate. We instead isolate a geopolitical signal by holding prompt content constant and varying only country name and language. The aggregate over three non-English languages is used as the reference rather than each country's official language, because tying the non-English side to each country's spoken language would leave most allied and neutral countries without a comparable score and break the balanced design. We name this asymmetry as a limitation later.

\section{Related Work}

\citet{barbieri-etal-2022-xlm} introduce the XLM-T family of Twitter-fine-tuned multilingual models, including the sentiment classifier we use. Their paper reports per-language accuracy on the UMSAB benchmark and establishes that XLM-R-Twitter generalizes across eight languages. Our work uses their model as an instrument and shifts the question from per-language accuracy to cross-lingual consistency on a fixed, controlled input.

\citet{faisal-anastasopoulos-2023-geographic} probe geographic and geopolitical biases in language models through factual question answering. They find that models under-represent the global South and encode viewpoints tied to dominant training languages. Our study extends this from factuality to sentiment and from open-ended queries to a balanced-design experiment, which lets us attach an effect size to the cross-lingual shift.

\citet{li-etal-2024-land} show that language models give language-dependent answers on contested territorial disputes, with the answer tracking the language's dominant political stance. Their finding that political framing leaks through the choice of prompt language is the phenomenon our sentiment Gap is built to measure at a continuous scale.

\citet{camara-etal-2022-mapping} measure intersectional demographic biases in multilingual sentiment systems across English, Spanish, and Arabic. They find Arabic-side degradation on sentiment fairness metrics. Their result grounds our inclusion of Arabic. Our target is geopolitical rather than demographic.

\citet{goldfarb-tarrant-etal-2023-cross} show that cross-lingual transfer amplifies gender and racial bias in sentiment classifiers compared to monolingual baselines. Their finding that transfer introduces directional bias supports the theoretical premise of our hypothesis. We differ in targeting country-category bias rather than demographic bias.

\section{Method}

\paragraph{Model.} We use twitter-xlm-roberta-base-sentiment \citep{barbieri-etal-2022-xlm}, trained on 198M tweets and fine-tuned on UMSAB for three-class sentiment. We chose it because it is the only widely available multilingual model that is both pretrained on large-scale social media across 100+ languages and ships with a fine-tuned sentiment head; training our own head on mBERT or mT5 would inject artifacts that confound the cross-lingual comparison.

\paragraph{Templates.} We wrote 30 English templates, organized into five structural families of six sentences each, with syntactic variety across the \texttt{\{country\}} slot's grammatical position. Families are News/media (``I saw a news report about \texttt{\{country\}}''), Education/reference (``We studied \texttt{\{country\}} in class''), Travel/location (``A colleague of mine is from \texttt{\{country\}}''), Conversation/mention (``\texttt{\{country\}} came up in our conversation''), and Descriptive/geographical (``\texttt{\{country\}} has a capital''). Templates were written to avoid mental-state verbs, evaluative adjectives, and presupposition-carrying constructions. Candidate templates that evoked a specific news event during informal review were rewritten. The full set, in all four languages, is in \texttt{main.ipynb}.

\paragraph{Countries.} We selected 30 countries partitioned into three balanced groups of 10. Adversaries are countries named as strategic competitors in the U.S.~National Security Strategy (2022), on the State Sponsors of Terrorism list, or under comprehensive U.S. sanctions: China, Russia, Iran, North Korea, Cuba, Syria, Venezuela, Belarus, Nicaragua, Myanmar. Allies are NATO member states with geographic spread: United Kingdom, France, Germany, Canada, Italy, Spain, Netherlands, Poland, Turkey, Norway. Neutral countries are declared-neutral or prominent Non-Aligned Movement members: Switzerland, Austria, Ireland, Serbia, India, Indonesia, Brazil, South Africa, Mexico, Argentina. Sweden and Finland are excluded because they joined NATO in 2023--2024.

The adversary/ally framing is U.S.-centric. The model under study is trained predominantly on English-language social media, which over-represents U.S. and allied media discourse, so a U.S.-centric geopolitical frame is the right frame against which to test whether the model's sentiment predictions track the framing of its dominant training distribution. We return to this asymmetry in Limitations.

\paragraph{Languages and translation.} The four languages are English, Chinese, Russian, Arabic. Chinese and Arabic translations were reviewed by the team member fluent in that language, checking semantic fidelity, absence of new sentiment-bearing vocabulary, and grammatical naturalness, including verb and adjective morphology that agrees with the \texttt{\{country\}} slot. No team member is a native Russian speaker, so Russian templates were validated by two-system MT consensus (Google Translate and DeepL) plus back-translation into English, with an additional gender-agreement check in predicate position. During validation we identified six templates (Arabic E1, E5, E6 and Russian D2, D6, E6) whose morphology agreed with the \texttt{\{country\}} slot and therefore produced a grammaticality gradient correlated with country-name gender. All six were rewritten to gender-invariant constructions (nominal sentences, impersonal plurals, reframings that move the grammatical subject to a gender-invariant noun). Russian validation is methodologically weaker than Chinese or Arabic validation, and we revisit this in Limitations.

\paragraph{Measurement.} The model returns a three-class softmax over \{Negative, Neutral, Positive\}. We read class-to-index mapping from \texttt{model.config.id2label} rather than hardcoding. The primary per-template score is bipolar valence, $S = P_{\text{pos}} - P_{\text{neg}}$, averaged over the 30 templates to give one $S_L(c)$ per country $c$ and language $L$. Bipolar valence is used rather than $P_{\text{neg}}$ alone because cross-lingual shifts can move probability mass between any two classes. $P_{\text{neg}}$ is retained as a sensitivity check.

The Cross-Lingual Sentiment Gap is defined as
\begin{equation}
  \text{Gap}(c) = S_{\text{EN}}(c) - \tfrac{1}{3}\bigl(S_{\text{ZH}}(c) + S_{\text{RU}}(c) + S_{\text{AR}}(c)\bigr),
\end{equation}
where a negative Gap means the model assigns more negative sentiment in English than in the non-English aggregate. Aggregating across ZH, RU, AR rather than tying the reference to each country's official language keeps every country on a comparable scale and preserves the balanced design. The tradeoff is that the aggregate does not correspond to any single country's actual spoken language, which we name as a limitation.

\paragraph{Statistical design.} The primary test is a one-way ANOVA on the 30 per-country Gap values across the three categories, combined with a directional contrast with weights $[+2, -1, -1]$ applied to $[\text{Adversary}, \text{Ally}, \text{Neutral}]$ means. The hypothesis is considered supported only if both the ANOVA clears $\alpha = .05$ and the contrast clears $\alpha = .05$ one-tailed in the predicted direction. We use Type II sums of squares. Assumption checks are Levene's test for variance homogeneity and Shapiro--Wilk on residuals. Because $S \in [-1, 1]$ is bounded, we also report a Kruskal--Wallis rank-based sensitivity check and a $P_{\text{neg}}$-only rerun of both the ANOVA and the contrast. Effect sizes are $\eta^2$ for the ANOVA and Cohen's d for the contrast, with pooled SD from residual MSE. Post-hoc pairwise comparisons use Holm--Bonferroni. A supplementary two-way ANOVA (language $\times$ category) on raw scores is reported as confirmation. All code is in \texttt{main.ipynb}.

\paragraph{Nonsense-country baseline.} We run the same 30 templates $\times$ 4 languages with the \texttt{\{country\}} slot replaced by five fictitious country names (Zorbia, Kelandor, Traviska, Nuberia, Sambalos) chosen to span the token-length distribution of real country names. The baseline gives a reference range for the Gap in the absence of country-identity signal. We use it descriptively rather than inferentially, because $n = 5$ does not support a formal test.

\section{Results}

\begin{table}[t]
  \centering
  \small
  \begin{tabular}{lccccc}
    \toprule
    Category & EN & ZH & RU & AR & Gap \\
    \midrule
    Adversary & $-$0.074 & $-$0.034 & $-$0.009 & $-$0.031 & $-$0.049 \\
    Ally      & $-$0.009 & $-$0.013 & $-$0.011 & $\phantom{-}$0.008 & $-$0.003 \\
    Neutral   & $-$0.002 & $-$0.009 & $-$0.012 & $\phantom{-}$0.008 & $\phantom{-}$0.002 \\
    \bottomrule
  \end{tabular}
  \caption{Category means of bipolar sentiment $S = P_{\text{pos}} - P_{\text{neg}}$ by language, and Cross-Lingual Sentiment Gap per category.}
  \label{tab:category-means}
\end{table}

\begin{figure}[t]
  \centering
  \includegraphics[width=\columnwidth]{figures/fig1_heatmap.pdf}
  \caption{Mean bipolar sentiment $S = P_{\text{pos}} - P_{\text{neg}}$ assigned by XLM-RoBERTa-Twitter to each country across four languages. Rows are grouped by category (adversary, ally, neutral, with a nonsense-country reference band).}
  \label{fig:heatmap}
\end{figure}

\begin{figure}[t]
  \centering
  \includegraphics[width=\columnwidth]{figures/fig2_interaction.pdf}
  \caption{Interaction plot of mean $S$ by language and category, with $\pm 1$ SE across the 10 countries per cell. Non-parallel lines indicate an interaction effect.}
  \label{fig:interaction}
\end{figure}

Table~\ref{tab:category-means} shows the category means. Adversaries score $-$0.049 on Gap, allies $-$0.003, and neutrals $+$0.002, so the English-minus-non-English disparity concentrates on adversaries. The one-way ANOVA on the 30 per-country Gap values across the three categories is significant, F(2, 27) = 8.75, p = .001, $\eta^2$ = .39. Levene's test for equal variances is not significant (p = .538) and Shapiro--Wilk on residuals is marginal (p = .071), so the OLS assumptions are acceptable.

The directional contrast with weights $[+2, -1, -1]$ confirms the predicted pattern: adversary Gap is more negative than the ally-plus-neutral mean, t(27) = $-$4.16, one-tailed p = .0001. The standardized mean difference between adversary (n=10) and pooled non-adversary (n=20) Gaps is Cohen's d = $-$1.61, large by conventional thresholds. Holm-adjusted pairwise comparisons show adversary vs ally p = .011 and adversary vs neutral p = .008, while ally vs neutral is null (p = .625). This means the hypothesis is supported for the adversary-versus-rest contrast but the ally/neutral distinction is not detectable at this sample size.

The Kruskal--Wallis rank-based check gives H = 11.07, p = .004, matching the parametric result. The $P_{\text{neg}}$-only sensitivity rerun gives F(2, 27) = 7.94, p = .002 on the ANOVA and t(27) = 3.95, one-tailed p = .0003 on the contrast, so the conclusion does not depend on the choice of bipolar valence over $P_{\text{neg}}$. The supplementary two-way ANOVA on raw scores shows a language $\times$ category interaction with partial $\eta^2 = .148$, consistent with Figure~\ref{fig:interaction}: adversary lines are pulled down in English relative to the other three languages while ally and neutral lines remain roughly parallel.

Figure~\ref{fig:heatmap} compares real-country Gaps to the nonsense-country reference band. The five nonsense countries produce Gaps in $[-0.068, -0.015]$ with mean $-$0.042, so the English-minus-non-English difference is modestly negative even when the country slot carries no identity. Adversary ($-$0.049) sits at this baseline while ally ($-$0.003) and neutral ($+$0.002) sit well above it. Adversary and nonsense Gaps therefore occupy the same range, while familiar ally and neutral names produce smaller Gaps. The baseline does not by itself identify which language side carries the asymmetry: a per-language look at the nonsense countries shows the negative Gap there is shaped partly by non-English scores on unfamiliar tokens, not the English side alone.

The per-language decomposition shows the aggregate Gap for adversaries is consistently negative across ZH, RU, and AR rather than driven by a single language. The effect is therefore a property of the English side being more negative, not of any one non-English language being more positive.

\section{Discussion}

Our results support the hypothesis of a Cross-Lingual Sentiment Gap. The one-way ANOVA across country categories on 30 per-country Gap values is significant (F(2, 27) = 8.75, p = .001, $\eta^2$ = .39), and the planned contrast confirms that the adversary Gap (M = $-$0.049) differs from the mean of ally and neutral Gaps (M $\approx$ 0) in the predicted direction (t(27) = $-$4.16, one-tailed p = .0001, d = $-$1.61). The nonsense-country baseline has mean Gap $-$0.042 in the range $[-0.068, -0.015]$. Adversary Gaps overlap this range while ally and neutral Gaps sit well above it, meaning familiar ally and neutral names produce smaller cross-lingual disparities than either adversary or fictitious names, whose disparities are of comparable magnitude. This is consistent with XLM-RoBERTa-Twitter encoding geopolitically-aligned sentiment differences that vary with the language of the prompt, even though the baseline is descriptive and does not on its own identify which language side carries the asymmetry.

The root-cause question has two plausible answers: the sentiment classifier head itself, or the underlying LM representations it reads from. Our design cannot fully separate these, but the fact that the effect is specific to adversaries and not to allies or neutrals, and that it survives a Kruskal--Wallis rank check and a $P_{\text{neg}}$-only rerun, rules out generic template valence as the driver. The most likely account is that English-heavy pretraining data associates adversary country names with more negative contexts than any of the three non-English corpora do, so the English representation of adversary country names enters the classifier already tilted. The classifier then amplifies that tilt. Both levels contribute, and a probing study on the frozen encoder is the natural next step.

The ally-versus-neutral null is also informative. It shows the effect is not a continuous ``how friendly is this country to the U.S.''~gradient but a categorical adversary-versus-rest tilt, which matches how the adversary label is most salient in U.S.~political discourse.

\section{Limitations}

We note the following limits. Country names in Russian and Arabic carry grammatical case and gender that do not uniformly agree with the verbs and adjectives in a fixed-string template. We rewrote six templates that were shown to produce a grammaticality gradient correlated with country-name gender, but residual case-mismatch noise in other Russian templates remains a limitation we do not quantify here. The single classifier (XLM-RoBERTa-Twitter) also limits generalization: whether other multilingual sentiment models show the same pattern is an open question. Template register is formal-to-neutral and findings may not transfer to conversational or social-media input. With 30 countries and four languages, statistical power is sufficient for category-level inference but not for per-country claims. The nonsense-country baseline is descriptive, not a formal control: $n = 5$ is too small for a hypothesis test, and fictitious names still tokenize to unfamiliar subwords, so the probe isolates country-identity signal imperfectly.

The language-country mapping is an oversimplification. Arabic is spoken in allied (Jordan), neutral, and adversary states, Chinese is spoken outside the PRC, and Russian is spoken well beyond Russia. The non-English aggregate is therefore a proxy rather than a direct mapping from language to a country's actual discourse environment. The adversary/ally partition is itself U.S.-centric, and the same experiment run with a different geopolitical frame would relabel the countries and likely produce a different Gap pattern. The translation protocol is also asymmetric: Chinese and Arabic were reviewed by fluent team members, while Russian relied on MT consensus and back-translation, and no formal inter-rater reliability was computed.

\section{Conclusion}

XLM-RoBERTa-Twitter assigns more negative sentiment to U.S.-designated adversary countries in English than in the non-English aggregate of Chinese, Russian, and Arabic, with a large effect size that survives rank-based and measure-based sensitivity checks. A nonsense-country baseline shows adversary Gaps overlap the range produced by country-identity-free names while ally and neutral Gaps sit above it, so the disparity is a specific adversary-versus-familiar signal rather than a continuous friend-foe gradient. Ally and neutral categories are indistinguishable from each other. Probing the frozen encoder to localize where this signal enters the pipeline is the natural follow-up.

\section*{Statement of Contributions}

The entire group contributed to researching topics for the paper. We worked individually to explore the literature on multilingual sentiment bias and to draft candidate templates and country categorizations before agreeing on the final design. We collaborated on all aspects of the project: template drafting, the inference pipeline with XLM-RoBERTa, the statistical analysis (one-way ANOVA, planned contrast, and sensitivity checks), visualizations, and the write-up. Translation validation was specialized by language competence: Nabil reviewed the Arabic translations for semantic fidelity, sentiment neutrality, and morphological agreement; Sebastien reviewed the Chinese translations under the same criteria. Russian validation, for which no team member is a native speaker, was conducted jointly by all three authors using the two-system MT consensus and back-translation procedure described in Method.

\bibliography{custom}

\end{document}
